\subsection{scanf / printf: форматы (C++17/20/23)}
\begin{itemize}[leftmargin=*, nosep]
\item Заголовок \verb|<cstdio>|. \verb|printf| получает значения, \verb|scanf| -- адреса: \verb|scanf("%d%lld", &n, &x);|. Для массива \verb|char s[100]| передавать \verb|s|, не \verb|&s|.
\item Тип обязан совпадать с форматом; иначе UB. При \verb|#define int long long| для таких переменных нужен \verb|%lld|, не \verb|%d|.
\end{itemize}

\subsection{Спецификаторы}
\begin{itemize}[leftmargin=*, nosep]
\item \verb|%d| -- \verb|int| (десятичный), \verb|%i| -- тот же вывод; при вводе \verb|%i| определяет основание: \verb|010| -- 8, \verb|0x10| -- 16. Для обычного ввода чисел брать \verb|%d|.
\item \verb|%u| -- \verb|unsigned int|; \verb|%o| -- восьмеричный; \verb|%x| / \verb|%X| -- шестнадцатеричный (на выводе буквы маленькие / большие). Для \verb|o, x, X| тоже unsigned-тип.
\item \verb|%ld| / \verb|%lu| -- \verb|long| / \verb|unsigned long|; \verb|%lld| / \verb|%llu| -- \verb|long long| / \verb|unsigned long long|. Размер \verb|long| зависит от платформы.
\item \textbf{Вещественные:} \verb|%f| -- обычная запись, \verb|%e| -- экспонента, \verb|%g| -- выбор между ними (убирает конечные нули дробной части), \verb|%a| -- шестнадцатеричная запись с двоичной экспонентой. Есть \verb|F, E, G, A| с заглавными буквами. Для scanf все эти формы принимают вещественный ввод, а не только соответствующий вид записи.
\item \textbf{printf:} \verb|%f| и \verb|%lf| получают \verb|double|; \verb|float| автоматически расширяется до \verb|double|. \verb|%Lf| -- \verb|long double|.
\item \textbf{scanf:} \verb|%f| требует \verb|float*|, \verb|%lf| -- \verb|double*|, \verb|%Lf| -- \verb|long double*|. Аналогично для \verb|e, g, a|.
\item \verb|%c| -- символ (printf получает \verb|int| после продвижения \verb|char|); \verb|%s| -- C-строка, при вводе до пробельного символа. \verb|%lc| / \verb|%ls| -- широкий символ / строка (\verb|wint_t| / \verb|wchar_t*| для printf; \verb|wchar_t*| для scanf).
\item \verb|%p| -- указатель: printf получает \verb|void*|, scanf -- \verb|void**|; для вывода указатель приводить к \verb|void*|. \verb|%%| -- сам знак процента, без аргумента.
\item \verb|%n| -- ничего не выводит / не читает; записывает число уже выведенных / прочитанных символов в \verb|int*|. Модификатор длины меняет тип указателя. У scanf не входит в число успешных присваиваний.
\end{itemize}

\subsection{printf: флаги, ширина, точность}
\begin{itemize}[leftmargin=*, nosep]
\item Формат: \verb|%[flags][width][.precision][length]specifier|.
\item Все стандартные флаги: \verb|-| -- выровнять влево; \verb|+| -- всегда показывать знак; пробел -- место вместо плюса; \verb|0| -- добивать число нулями слева после знака/префикса; \verb|#| -- альтернативная форма. \verb|+| сильнее пробела, \verb|-| отключает \verb|0|. Флаги применимы не ко всем спецификаторам.
\item \verb|#|: для \verb|o| гарантирует начальный 0; для ненулевых \verb|x/X| добавляет \verb|0x/0X|; для вещественных сохраняет точку, а у \verb|g/G| ещё и конечные нули. Для \verb|d, u, s, c| не применять.
\item Ширина -- \textbf{минимум} символов, ничего не обрезает: \verb|%6d|; \verb|%06d| с числом 42 даёт \verb|000042|; \verb|%-6d| добавляет пробелы справа.
\item Точность: \verb|%.3f| -- 3 цифры после точки (с округлением); \verb|%.3e| -- 3 после точки в мантиссе; \verb|%.3g| -- 3 значащие цифры. По умолчанию у \verb|f/e/g| точность 6; у \verb|g| точность 0 трактуется как 1.
\item Для целых \verb|%.4d| -- минимум 4 цифры; заданная точность отключает флаг \verb|0|. Ноль с точностью 0 даёт пустой вывод (исключение: \verb|%#.0o| даёт 0). Для строк \verb|%.5s| -- максимум 5 байтов.
\item \verb|*| берёт ширину / точность из аргумента типа \verb|int|: \verb|printf("%*.*f", w, p, x);|. Отрицательная ширина означает \verb|-| и её модуль; отрицательная точность считается неуказанной. Даже при \verb|#define int long long| для звёздочек передавать именно \verb|(int32_t)w| на обычной 32-битной модели int.
\end{itemize}

\subsection{Модификаторы длины}
\begin{itemize}[leftmargin=*, nosep]
\item Для целых: \verb|hh| -- \verb|signed/unsigned char|; \verb|h| -- \verb|short|; \verb|l| -- \verb|long|; \verb|ll| -- \verb|long long|; \verb|j| -- \verb|intmax_t/uintmax_t|; \verb|z| -- \verb|size_t| для unsigned, его знаковый аналог для signed; \verb|t| -- \verb|ptrdiff_t| для signed, его unsigned-аналог для unsigned.
\item Примеры: \verb|%hhd|, \verb|%hu|, \verb|%zu| для \verb|size_t|, \verb|%td| для \verb|ptrdiff_t|. В scanf нужны указатели на точные типы; в printf узкие целые проходят стандартные продвижения. \verb|L| -- \verb|long double|; \verb|l| также используется с вещественными, \verb|c/s|, как описано выше. Стандартного формата для \verb|__int128| нет.
\end{itemize}

\subsection{scanf: ширина, пропуск и пробелы}
\begin{itemize}[leftmargin=*, nosep]
\item Формат: \verb|%[*][width][length]specifier|. Здесь нет флагов printf и точности. \verb|*| \textbf{подавляет присваивание}: \verb|scanf("%*d%d", &x);| пропускает первое число. Аргумент под пропускаемое поле не нужен.
\item Ширина -- \textbf{максимум} прочитанных символов поля: \verb|char s[100]; scanf("%99s", s);|. Для \verb|%s| и \verb|%[| нужен ещё байт под \verb|'\0'|. \verb|%3c| читает 3 символа в буфер и не добавляет \verb|'\0'|.
\item \verb|%[abc]| -- непустая последовательность символов из набора; \verb|%[^\n]| -- до перевода строки. \verb|]| включается первым в наборе (после возможного \verb|^|); \verb|-| надёжнее ставить первым / последним. Диапазоны вроде \verb|[a-z]| поддерживаются обычными библиотеками, но зависят от реализации.
\item Почти все преобразования пропускают начальные пробелы; исключения -- \verb|%c|, \verb|%[|, \verb|%n|. Поэтому \verb|scanf(" %c", &c);| читает следующий непробельный символ.
\item Любой пробельный символ в формате, включая \verb|\n|, поглощает произвольное число пробельных символов ввода. Не ставить пробел / \verb|\n| в конце формата интерактивного ввода: вызов может ждать следующий непробельный символ. Обычные символы формата должны совпасть со входом: \verb|"%d,%d"| требует запятую.
\item \verb|scanf| возвращает число успешно присвоенных полей; \verb|EOF| -- если ввод закончился / сломался до первого преобразования. Читать до конца: \verb|while (scanf("%d", &x) == 1)|. \verb|printf| возвращает число выведенных символов либо отрицательное значение при ошибке.
\item После \verb|ios::sync_with_stdio(false)| не смешивать эти функции с \verb|cin/cout|. \verb|%b|, \verb|%m|, \verb|%I64d| и другие расширения не считать переносимыми для C++17/20/23.
\end{itemize}
\par \textbf{Источник:} \href{https://www.open-std.org/jtc1/sc22/wg14/www/docs/n1570.pdf}{C11, 7.21.6.1--7.21.6.4 (printf/scanf)}.
